 



























\documentclass[11pt,a4paper]{article}

% ============================================================
% Packages et mise en page
% ============================================================
\usepackage[utf8]{inputenc}
\usepackage[T1]{fontenc}
\usepackage[french]{babel}
\usepackage{lmodern}
\usepackage{microtype}
\usepackage{geometry}
\geometry{margin=2.0cm}
\usepackage{amsmath,amssymb,mathtools}
\usepackage{booktabs,longtable,tabularx,array}
\usepackage{enumitem}
\usepackage{float}
\usepackage{tikz}
\usetikzlibrary{arrows.meta,positioning,shapes.geometric,fit}
\usepackage{adjustbox}
\usepackage{xcolor}
\usepackage{hyperref}

\setlength{\parindent}{0pt}
\setlength{\parskip}{5pt}
\setlist{nosep,leftmargin=*}
\hypersetup{colorlinks=true,linkcolor=blue,urlcolor=blue}

% Graphes contenus dans la largeur du texte
\tikzset{
  nodebox/.style={
    draw,rounded corners,align=center,
    text width=2.7cm,minimum height=0.8cm,
    font=\small
  },
  smallbox/.style={
    draw,rounded corners,align=center,
    text width=2.2cm,minimum height=0.7cm,
    font=\footnotesize
  },
  arrow/.style={-{Stealth[length=2mm]},thick},
  dashedarrow/.style={-{Stealth[length=2mm]},thick,dashed}
}

\title{\textbf{EX12 -- Organisation adaptative des centres AdaptiveRBF}\\
\large Sélection progressive, guidage des centres faibles et empreintes relationnelles}
\author{Protocole expérimental}
\date{Version 1.0 -- Spécification pré-enregistrable}

\begin{document}
\maketitle

\begin{abstract}
Ce protocole étudie si une organisation explicite des centres d'un réseau
AdaptiveRBF améliore la précision, la stabilité et l'efficacité de
l'apprentissage. Trois mécanismes sont séparés :
(i) sélection progressive des centres selon leur contribution prédictive ;
(ii) guidage des centres peu utiles par les centres utiles et par les
résidus du modèle ;
(iii) organisation relationnelle des centres à partir d'une empreinte
calculée sur leur position, leur activation, leur contribution et leur
redondance.

Les mécanismes sont évalués séparément puis combinés. L'objectif n'est
pas de présumer leur efficacité, mais de déterminer dans quelles conditions
ils améliorent ou dégradent AdaptiveRBF.
\end{abstract}

\tableofcontents
\newpage

% ============================================================
\section{Question scientifique}
% ============================================================

\textbf{Question principale :}

Peut-on améliorer un réseau AdaptiveRBF en apprenant explicitement
l'utilité, la spécialisation et les relations entre ses centres,
plutôt que de laisser tous les centres être optimisés uniquement
par la perte globale ?

Les hypothèses sont :

\begin{itemize}
  \item H1 : la contribution prédictive d'un centre est un meilleur
  critère de sélection que sa seule activation moyenne.
  \item H2 : une sélection progressive avec réentraînement peut
  réduire le nombre de centres sans perte importante de précision.
  \item H3 : les centres faibles peuvent devenir utiles s'ils sont
  guidés vers les régions où le modèle commet encore des erreurs.
  \item H4 : une empreinte multidimensionnelle permet d'identifier
  les centres redondants, sous-utilisés ou complémentaires.
  \item H5 : une organisation relationnelle des centres améliore
  la stabilité entre graines et la généralisation.
  \item H6 : les bénéfices éventuels persistent face à des variantes
  AdaptiveRBF de capacité et de budget comparables.
\end{itemize}

Toutes les hypothèses sont falsifiables. Une amélioration sur une seule
fonction ou une seule graine ne suffit pas à confirmer une hypothèse.

% ============================================================
\section{Architecture de référence : AdaptiveRBF classique}
% ============================================================

\subsection{Notation}

Soit un jeu d'entraînement :
\[
\mathcal{D}_{\mathrm{train}}
=\{(x_n,y_n)\}_{n=1}^{N},
\quad x_n\in\mathbb{R}^{p},
\quad y_n\in\mathbb{R}^{q}.
\]

L'encodeur transforme l'entrée :
\[
z_n=\tanh(W_e x_n+b_e),
\qquad z_n\in\mathbb{R}^{d}.
\]

Le réseau possède \(K\) centres :
\[
C=\{c_1,\ldots,c_K\},\qquad c_k\in\mathbb{R}^{d}.
\]

Chaque centre possède une largeur positive :
\[
\sigma_k=\sigma_{\min}
+(\sigma_{\max}-\sigma_{\min})\operatorname{sigmoid}(s_k).
\]

L'activation du centre \(k\) pour l'entrée \(x_n\) est :
\[
\phi_{nk}
=
\exp\left(
-\frac{\|z_n-c_k\|_2^2}{2\sigma_k^2}
\right).
\]

Les activations sont assemblées :
\[
\phi_n=[\phi_{n1},\ldots,\phi_{nK}].
\]

La sortie est :
\[
\hat y_n=W_o\phi_n+b_o.
\]

La perte de référence est :
\[
\mathcal{L}_{\mathrm{task}}
=
\frac{1}{N}\sum_{n=1}^{N}
\ell(\hat y_n,y_n).
\]

Pour une régression, la métrique d'entraînement principale est la MSE.
La MAE et la RMSE sont également rapportées.

\subsection{Initialisation et entraînement de référence}

\begin{enumerate}
  \item Ajuster les éventuelles normalisations uniquement sur
  \(\mathcal{D}_{\mathrm{train}}\).
  \item Entraîner l'encodeur de référence ou l'initialiser selon
  la configuration pré-enregistrée.
  \item Initialiser les centres par KMeans sur les représentations
  latentes des exemples d'entraînement uniquement.
  \item Initialiser les largeurs à partir des distances entre centres,
  puis les rendre apprenables.
  \item Entraîner encodeur, centres, largeurs et readout avec
  l'optimiseur et le budget fixés.
  \item Sélectionner le checkpoint uniquement sur validation.
\end{enumerate}

\subsection{Graphe de référence}

\begin{figure}[H]
\centering
\begin{adjustbox}{max width=\linewidth}
\begin{tikzpicture}[node distance=0.65cm]
\node[nodebox] (x) {Entrée \(x\)};
\node[nodebox,right=of x] (enc) {Encodeur\\\(z=\tanh(W_ex+b_e)\)};
\node[nodebox,right=of enc] (rbf) {Centres RBF\\\(\phi_k(z,c_k,\sigma_k)\)};
\node[nodebox,right=of rbf] (out) {Readout linéaire\\\(\hat y=W_o\phi+b_o\)};
\node[nodebox,right=of out] (y) {Prédiction \(\hat y\)};
\draw[arrow] (x)--(enc);
\draw[arrow] (enc)--(rbf);
\draw[arrow] (rbf)--(out);
\draw[arrow] (out)--(y);
\end{tikzpicture}
\end{adjustbox}
\caption{AdaptiveRBF classique : les centres sont optimisés par la perte
globale, sans mécanisme explicite d'organisation entre centres.}
\end{figure}

% ============================================================
\section{Mesurer la force réelle d'un centre}
% ============================================================

\subsection{Pourquoi l'activation seule ne suffit pas}

Un centre peut avoir une activation élevée sans être indispensable :
un autre centre peut fournir une réponse similaire.

Inversement, un centre rarement activé peut être essentiel dans une
région rare de l'espace d'entrée.

La force d'un centre ne sera donc pas définie uniquement par
\(\mathbb{E}[\phi_k]\). Elle sera mesurée par plusieurs indicateurs.

\subsection{Activation moyenne}

\[
A_k=\frac{1}{N}\sum_{n=1}^{N}\phi_{nk}.
\]

Cet indicateur mesure la fréquence et l'intensité d'activation.
Il ne mesure pas directement l'utilité prédictive.

\subsection{Contribution par ablation}

Soit \(\mathcal{L}_{\mathrm{full}}\) la perte du modèle complet.
On retire le centre \(k\), on réajuste le readout selon une procédure
fixée, puis on mesure la perte \(\mathcal{L}_{-k}\).

\[
\Delta_k^{\mathrm{abl}}
=
\mathcal{L}_{-k}-\mathcal{L}_{\mathrm{full}}.
\]

Interprétation :
\begin{itemize}
  \item \(\Delta_k^{\mathrm{abl}}>0\) : le centre apporte une contribution utile ;
  \item \(\Delta_k^{\mathrm{abl}}\approx0\) : le centre est possiblement redondant ;
  \item \(\Delta_k^{\mathrm{abl}}<0\) : le centre nuit à la solution actuelle.
\end{itemize}

L'ablation doit être réalisée sur validation, avec une procédure identique
pour tous les centres. Le test scellé ne sert jamais à classer les centres.

\subsection{Contribution résiduelle}

Le résidu du modèle est :
\[
r_n=y_n-\hat y_n.
\]

Pour mesurer si un centre explique une partie du résidu, on considère
une mise à jour locale du readout utilisant uniquement son activation.
La réduction de perte obtenue est notée :
\[
\Delta_k^{\mathrm{res}}.
\]

Cette mesure estime si le centre peut corriger une erreur que le modèle
commet encore.

\subsection{Redondance entre centres}

La similarité d'activation entre deux centres est :
\[
R_{kj}
=
\frac{
\sum_n(\phi_{nk}-\bar\phi_k)(\phi_{nj}-\bar\phi_j)
}{
\sqrt{\sum_n(\phi_{nk}-\bar\phi_k)^2}
\sqrt{\sum_n(\phi_{nj}-\bar\phi_j)^2}
+\epsilon
}.
\]

Une valeur proche de \(1\) indique des profils d'activation similaires.
La redondance doit être interprétée avec la contribution : deux centres
similaires peuvent néanmoins avoir des poids de sortie différents.

% ============================================================
\section{Empreinte d'un centre}
% ============================================================

\subsection{Définition}

L'empreinte du centre \(k\) est un vecteur descriptif :
\[
e_k=
[
\tilde c_k,\,
\log\sigma_k,\,
A_k,\,
\Delta_k^{\mathrm{abl}},\,
\Delta_k^{\mathrm{res}},\,
v_k,\,
u_k
].
\]

Où :

\begin{itemize}
  \item \(\tilde c_k\) : position normalisée du centre dans l'espace latent ;
  \item \(\log\sigma_k\) : largeur du centre ;
  \item \(A_k\) : activation moyenne ;
  \item \(\Delta_k^{\mathrm{abl}}\) : contribution mesurée par ablation ;
  \item \(\Delta_k^{\mathrm{res}}\) : capacité à réduire le résidu ;
  \item \(v_k\) : variance ou dispersion de ses activations ;
  \item \(u_k\) : mesure d'utilisation sur les données.
\end{itemize}

Les composantes sont normalisées avec des statistiques calculées
uniquement sur le train ou la validation autorisée. Les valeurs de
contribution sont estimées sur validation.

\subsection{Graphe des empreintes}

Chaque centre devient un nœud. Une arête relie deux centres si leurs
empreintes sont proches ou si leurs profils d'activation sont similaires.

\[
d_{kj}^{e}=\|e_k-e_j\|_2.
\]

\[
G_C=(V_C,E_C),\qquad
V_C=\{1,\ldots,K\}.
\]

Une arête est créée si :
\[
(k,j)\in E_C
\quad\Longleftrightarrow\quad
d_{kj}^{e}<\tau_e
\]
ou si \(j\) appartient aux \(m\) plus proches voisins de \(k\).

La règle retenue, le seuil \(\tau_e\) ou le nombre de voisins \(m\),
doivent être fixés sur validation et non sur test.

\begin{figure}[H]
\centering
\begin{adjustbox}{max width=0.85\linewidth}
\begin{tikzpicture}[node distance=1.1cm]
\node[smallbox] (c1) {Centre 1\\empreinte \(e_1\)};
\node[smallbox,right=of c1] (c2) {Centre 2\\empreinte \(e_2\)};
\node[smallbox,below=of c1] (c3) {Centre 3\\empreinte \(e_3\)};
\node[smallbox,below=of c2] (c4) {Centre 4\\empreinte \(e_4\)};
\draw[arrow] (c1)--(c2);
\draw[arrow] (c1)--(c3);
\draw[arrow] (c2)--(c4);
\draw[dashedarrow] (c3)--(c4);
\end{tikzpicture}
\end{adjustbox}
\caption{Graphe relationnel des centres. Les arêtes décrivent des
similarités ou relations apprises ; elles ne remplacent pas les
activations RBF.}
\end{figure}

\subsection{Mise à jour des empreintes}

Les empreintes sont recalculées périodiquement, et non à chaque exemple,
pour limiter le coût et le bruit.

À l'époque \(t\) :
\[
e_k^{(t)}=\operatorname{Fingerprint}(c_k^{(t)},\sigma_k^{(t)},
\Phi^{(t)},r^{(t)},W_o^{(t)}).
\]

Le graphe relationnel est ensuite reconstruit ou mis à jour selon
une cadence pré-enregistrée.

Important : les empreintes sont des statistiques descriptives.
Elles ne doivent pas utiliser les labels du test ni les prédictions
du test pour orienter l'apprentissage.

% ============================================================
\section{Variantes expérimentales}
% ============================================================

\subsection{B0 -- AdaptiveRBF classique}

Architecture et entraînement de référence.
Aucun classement, guidage ou graphe d'empreintes.

\subsection{B1 -- Classement des centres}

Les centres sont classés selon :
\[
S_k=
\alpha\widehat{\Delta}_k^{\mathrm{abl}}
+\beta\widehat{\Delta}_k^{\mathrm{res}}
-\gamma\widehat{\mathrm{Redundancy}}_k.
\]

Les coefficients sont fixés avant la campagne confirmatoire.
Le classement est rapporté, mais aucun centre n'est supprimé.

Objectif : vérifier si le classement identifie des centres utiles
et si cette utilité est stable entre graines.

\subsection{B2 -- Sélection progressive avec réentraînement}

Procédure :

\begin{enumerate}
  \item Entraîner le modèle complet.
  \item Mesurer \(S_k\) pour tous les centres.
  \item Retirer un petit groupe de centres de score faible.
  \item Réentraîner le modèle restant avec le même budget défini.
  \item Recalculer les contributions dans le nouveau modèle.
  \item Répéter jusqu'à atteindre le budget minimal ou un critère d'arrêt.
\end{enumerate}

Le retrait est progressif, car la contribution d'un centre dépend
des autres centres présents.

Critère d'arrêt possible :
\[
\mathcal{L}_{\mathrm{val}}^{(t)}
>
\mathcal{L}_{\mathrm{val}}^{(0)}+\delta
\]
ou lorsque le nombre minimal de centres est atteint.

Le seuil \(\delta\) est fixé avant les expériences.

\subsection{B3 -- Guidage des centres faibles par les centres forts}

Les centres forts sont définis par leur contribution prédictive,
pas seulement par leur activation.

On sépare :
\[
\mathcal{S}=\{k:S_k\geq\tau_S\},
\qquad
\mathcal{W}=\{k:S_k<\tau_S\}.
\]

Les centres faibles sont guidés vers les régions où le modèle commet
encore des erreurs, tout en utilisant les centres forts comme référence.

On définit un score d'erreur locale :
\[
q_n= \|r_n\|_2.
\]

La distribution de responsabilité cible d'un centre faible \(k\) est :
\[
p_{nk}
=
\frac{
\phi_{nk}\,q_n
}{
\sum_{m\in\mathcal{W}}\phi_{nm}q_n+\epsilon
}.
\]

Cette distribution encourage les centres faibles à se concentrer
sur les exemples mal expliqués, sans copier directement les centres forts.

Une perte de spécialisation possible est :
\[
\mathcal{L}_{\mathrm{special}}
=
-\frac{1}{N}\sum_n
q_n\log\left(
\sum_{k\in\mathcal{W}}\phi_{nk}+\epsilon
\right).
\]

Une perte de séparation évite que les centres faibles se rapprochent
tous du même centre fort :
\[
\mathcal{L}_{\mathrm{sep}}
=
\frac{1}{|\mathcal{W}||\mathcal{S}|}
\sum_{k\in\mathcal{W}}\sum_{j\in\mathcal{S}}
\exp\left(-\frac{\|c_k-c_j\|_2^2}{2\rho^2}\right).
\]

La perte totale devient :
\[
\mathcal{L}_{B3}
=
\mathcal{L}_{\mathrm{task}}
+\lambda_s\mathcal{L}_{\mathrm{special}}
+\lambda_d\mathcal{L}_{\mathrm{sep}}.
\]

Les coefficients \(\lambda_s,\lambda_d,\rho\) sont sélectionnés sur
validation et verrouillés avant le test.

\textbf{Précaution :} une forte activation d'un centre fort ne doit pas
être utilisée comme cible à copier. Le guidage vise la complémentarité
et les zones mal expliquées.

\subsection{B4 -- Organisation par empreintes}

B4 ajoute un module qui traite les empreintes des centres et produit
des corrections de position ou de largeur.

\[
\delta c_k,\delta s_k
=
F_\psi(e_k,\mathcal{N}(k)),
\]

où \(\mathcal{N}(k)\) est le voisinage du centre \(k\) dans le graphe
relationnel.

Les paramètres effectifs deviennent :
\[
c_k^{\mathrm{eff}}=c_k+\eta_c\delta c_k,
\]
\[
s_k^{\mathrm{eff}}=s_k+\eta_s\delta s_k,
\qquad
\sigma_k^{\mathrm{eff}}
=\sigma_{\min}
+(\sigma_{\max}-\sigma_{\min})
\operatorname{sigmoid}(s_k^{\mathrm{eff}}).
\]

Le module \(F_\psi\) peut être un petit réseau partagé appliqué
à chaque centre. Le partage limite le nombre de paramètres.

\subsection{Graphe complet de B4}

\begin{figure}[H]
\centering
\begin{adjustbox}{max width=\linewidth}
\begin{tikzpicture}[node distance=0.55cm]
\node[nodebox] (x) {Entrée \(x\)};
\node[nodebox,right=of x] (enc) {Encodeur\\latent \(z\)};
\node[nodebox,right=of enc] (rbf) {Centres RBF\\activations \(\phi_k\)};
\node[nodebox,right=of rbf] (out) {Readout\\prédiction \(\hat y\)};
\draw[arrow] (x)--(enc);
\draw[arrow] (enc)--(rbf);
\draw[arrow] (rbf)--(out);

\node[nodebox,below=1.0cm of rbf] (fp) {Empreintes\\\(e_k\)};
\node[nodebox,left=of fp] (stats) {Statistiques\\activation, erreur,\\contribution};
\node[nodebox,right=of fp] (graph) {Graphe des\\centres};
\node[nodebox,right=of graph] (org) {Organisateur\\\(F_\psi\)};
\draw[arrow] (rbf.south)--(fp.north);
\draw[arrow] (stats)--(fp);
\draw[arrow] (fp)--(graph);
\draw[arrow] (graph)--(org);
\draw[arrow] (org.north) |- (rbf.south);
\end{tikzpicture}
\end{adjustbox}
\caption{B4 : les empreintes décrivent les centres ; le graphe exprime
leurs relations ; l'organisateur ajuste leur configuration.}
\end{figure}

\textbf{Contraintes :}
\begin{itemize}
  \item l'organisateur ne reçoit pas les labels du test ;
  \item les empreintes sont calculées sur train et validation uniquement ;
  \item la mise à jour des empreintes est périodique et déterministe ;
  \item les corrections sont bornées pour éviter des déplacements instables ;
  \item le graphe relationnel est distinct du graphe CNC des opérations.
\end{itemize}

\subsection{B5 -- Variante combinée}

B5 combine B2, B3 et B4 :
\begin{itemize}
  \item sélection progressive ;
  \item guidage des centres faibles ;
  \item organisation relationnelle par empreintes.
\end{itemize}

B5 est évaluée seulement après l'analyse séparée de B2, B3 et B4.
Elle ne remplace pas les expériences isolées, car une amélioration
combinée ne permettrait pas d'identifier le mécanisme responsable.

% ============================================================
\section{Protocole d'entraînement}
% ============================================================

\subsection{Phases}

\begin{figure}[H]
\centering
\begin{adjustbox}{max width=0.88\linewidth}
\begin{tikzpicture}[node distance=0.55cm]
\node[nodebox] (d) {Données train\\validation\\test scellé};
\node[nodebox,below=of d] (base) {Entraîner\\AdaptiveRBF B0};
\node[nodebox,below=of base] (measure) {Mesurer les centres\\et calculer empreintes};
\node[nodebox,below=of measure] (variant) {Appliquer B1/B2/B3/B4};
\node[nodebox,below=of variant] (retrain) {Réentraînement\\selon la variante};
\node[nodebox,below=of retrain] (select) {Sélection checkpoint\\sur validation};
\node[nodebox,below=of select] (test) {Évaluation finale\\sur test scellé};
\draw[arrow] (d)--(base);
\draw[arrow] (base)--(measure);
\draw[arrow] (measure)--(variant);
\draw[arrow] (variant)--(retrain);
\draw[arrow] (retrain)--(select);
\draw[arrow] (select)--(test);
\end{tikzpicture}
\end{adjustbox}
\caption{Pipeline expérimental. Le test n'intervient qu'à l'évaluation finale.}
\end{figure}

\subsection{Budget et comparateurs}

Chaque variante doit être comparée à :

\begin{itemize}
  \item B0 AdaptiveRBF classique avec le même nombre initial de centres ;
  \item AdaptiveRBF avec nombre de centres réduit, mais sans sélection ;
  \item MLP direct de capacité comparable ;
  \item MLP plus large, comme référence de capacité ;
  \item variante sans empreintes, mais avec le même coût d'entraînement
  lorsque cela est possible.
\end{itemize}

Les comparaisons doivent rapporter à la fois :
(i) résultats à paramètres comparables et
(ii) résultats à budget de calcul comparable.

\subsection{Graines}

Utiliser au minimum dix graines pré-enregistrées :
\[
\{1501,1607,1717,2029,3037,4099,4517,5003,6011,7013\}.
\]

La graine contrôle :
\begin{itemize}
  \item l'échantillonnage des données synthétiques ;
  \item l'initialisation des poids ;
  \item l'initialisation KMeans ;
  \item l'ordre des mini-batches.
\end{itemize}

Les variantes d'une même graine utilisent les mêmes données et les
mêmes splits. Les comparaisons sont appariées.

\subsection{Données}

Utiliser deux familles de tâches.

\textbf{Famille A -- fonctions synthétiques :}
\[
f_1(x)=\sin(x),
\]
\[
f_2(x_1,x_2)=x_1x_2,
\]
\[
f_3(x_1,x_2)=x_1^2+x_2,
\]
\[
f_4(x_1,x_2,x_3)=\sin(x_1x_2)+x_3.
\]

\textbf{Famille B -- fonctions supplémentaires :}
\[
f_5(x_1,x_2)=\sin(x_1)+0.5\cos(2x_2),
\]
\[
f_6(x_1,x_2,x_3)=x_1x_2+0.25x_3^2.
\]

Les domaines, tailles train/validation/test, distributions et
règles d'extrapolation doivent être définis dans les configurations
versionnées avant lancement.

\textbf{Famille C -- données tabulaires :}

Utiliser les mêmes splits contrôlés que les expériences de référence.
Les données tabulaires servent à vérifier que les résultats ne sont
pas limités aux fonctions synthétiques.

\subsection{Optimisation}

Configuration initiale recommandée :
\begin{itemize}
  \item optimiseur : Adam ;
  \item learning rate : \(10^{-3}\) ;
  \item weight decay : \(10^{-5}\) ;
  \item batch size : 128 ;
  \item maximum : 800 époques ;
  \item early stopping : patience 50 ;
  \item checkpoint : meilleure MAE validation ;
  \item normalisation des cibles : train uniquement.
\end{itemize}

Toute modification de budget doit être appliquée de façon équitable
aux variantes comparées et documentée.

% ============================================================
\section{Métriques}
% ============================================================

\subsection{Précision}

Rapporter :
\[
\mathrm{MAE}=\frac{1}{N}\sum_n|y_n-\hat y_n|,
\]
\[
\mathrm{MSE}=\frac{1}{N}\sum_n(y_n-\hat y_n)^2,
\]
\[
\mathrm{RMSE}=\sqrt{\mathrm{MSE}}.
\]

\subsection{Stabilité}

Pour chaque variante, rapporter :
\begin{itemize}
  \item moyenne et écart-type entre graines ;
  \item médiane et intervalle interquartile ;
  \item pire résultat ;
  \item taux de runs dépassant le seuil d'erreur pré-enregistré.
\end{itemize}

\subsection{Structure des centres}

Rapporter :
\begin{itemize}
  \item distribution des scores \(S_k\) ;
  \item nombre de centres retenus ;
  \item taux de centres redondants ;
  \item diversité des profils d'activation ;
  \item évolution des empreintes au cours de l'entraînement ;
  \item stabilité du classement entre graines ;
  \item proportion de centres qui changent de catégorie.
\end{itemize}

\subsection{Coût}

Rapporter :
\begin{itemize}
  \item nombre de paramètres ;
  \item taille du checkpoint ;
  \item temps d'entraînement ;
  \item latence p50, p95 et p99 ;
  \item mémoire maximale si mesurée de manière fiable ;
  \item coût supplémentaire de calcul des empreintes ;
  \item coût de reconstruction du graphe relationnel.
\end{itemize}

Aucune revendication énergétique n'est permise sans mesure directe
de consommation.

% ============================================================
\section{Analyse statistique}
% ============================================================

Pour chaque tâche et chaque graine, calculer la différence appariée :
\[
d_s=\mathrm{MAE}_{\mathrm{variant},s}
-\mathrm{MAE}_{B0,s}.
\]

Une valeur négative favorise la variante.

Rapporter :
\begin{itemize}
  \item moyenne des différences appariées ;
  \item intervalle de confiance bootstrap à 95\% sur les graines ;
  \item taille d'effet ;
  \item nombre de graines gagnées, perdues et ex æquo.
\end{itemize}

Les graines, et non les lignes du test, constituent l'unité de répétition
pour l'inférence entre entraînements.

Les résultats par fonction sont rapportés séparément. Une moyenne globale
ne doit pas masquer des gains sur certaines fonctions et des pertes
sur d'autres.

Les comparaisons multiples entre B1--B5 sont corrigées selon une méthode
pré-enregistrée, par exemple Holm.

% ============================================================
\section{Ablations obligatoires}
% ============================================================

\begin{enumerate}
  \item Empreinte sans contribution par ablation.
  \item Empreinte sans résidu.
  \item Graphe relationnel sans correction des centres.
  \item Correction des centres sans graphe relationnel.
  \item Guidage sans terme de séparation.
  \item Sélection progressive sans réentraînement.
  \item Même nombre de centres retirés au hasard.
  \item Centres faibles guidés par les résidus, sans référence aux forts.
  \item B4 avec graphe aléatoire de même densité.
\end{enumerate}

Ces ablations permettent de vérifier si l'amélioration vient réellement
des empreintes et des relations entre centres, plutôt que d'une simple
augmentation de capacité ou d'un entraînement supplémentaire.

% ============================================================
\section{Risques et contrôles}
% ============================================================

\begin{longtable}{p{0.27\linewidth}p{0.63\linewidth}}
\toprule
\textbf{Risque} & \textbf{Contrôle} \\
\midrule
Activation forte confondue avec utilité &
Mesurer contribution par ablation et réduction du résidu. \\
Centres faibles supprimés trop tôt &
Réévaluer les scores après chaque réentraînement. \\
Centres faibles copient les forts &
Utiliser résidus et terme de séparation. \\
Empreintes instables &
Lissage temporel et cadence de mise à jour pré-enregistrée. \\
Fuite de données &
Aucune statistique de centre calculée sur le test. \\
Surapprentissage validation &
Limiter les choix d'hyperparamètres et conserver le test scellé. \\
Gain dû au budget supplémentaire &
Comparer à budget de calcul équivalent. \\
Graphe relationnel arbitraire &
Comparer au graphe aléatoire et aux ablations. \\
Résultat dépendant d'une graine &
Au moins dix graines appariées. \\
\bottomrule
\end{longtable}

% ============================================================
\section{Critères de conclusion}
% ============================================================

\textbf{Sélection progressive soutenue} si elle réduit le nombre de
centres tout en conservant une précision comparable et une stabilité
au moins équivalente à B0.

\textbf{Guidage soutenu} si B3 améliore les erreurs sur plusieurs tâches
et graines, et si les centres faibles deviennent plus complémentaires
plutôt que simplement plus activés.

\textbf{Empreintes soutenues} si B4 améliore de manière reproductible
la précision, la stabilité ou le compromis précision/coût, et si les
ablations montrent que le graphe relationnel apporte une contribution
spécifique.

\textbf{Hypothèse rejetée ou non soutenue} si les gains sont limités à
une tâche, disparaissent sur plusieurs graines, ou sont expliqués par
un budget supérieur.

Une variante plus précise mais beaucoup plus coûteuse n'est pas
automatiquement considérée comme meilleure. Le compromis doit être
rapporté explicitement.

% ============================================================
\section{Livrables et reproductibilité}
% ============================================================

Structure recommandée :
\begin{verbatim}
EX12/
  README_EX12.md
  protocol/
    EX12_protocol.tex
    preregistration.yaml
  configs/
    datasets/
    variants/
    optimizer/
  src/
    adaptive_rbf.py
    center_scores.py
    progressive_selection.py
    weak_center_guidance.py
    center_fingerprints.py
    center_graph.py
    fingerprint_organizer.py
  tests/
  runs/
    <task>/<variant>/<seed>/
      config.yaml
      checkpoint.pt
      predictions_val.csv
      metrics.json
      center_fingerprints.csv
      center_graph.json
      training_log.csv
  reports/
    aggregate_metrics.csv
    seed_stability.csv
    center_analysis/
    figures/
  manifests/
    checksums.sha256
\end{verbatim}

Chaque run doit enregistrer :
\begin{itemize}
  \item hash du code et de la configuration ;
  \item graine et versions logicielles ;
  \item checkpoint sélectionné ;
  \item prédictions train/validation ;
  \item métriques ;
  \item empreintes et graphe des centres ;
  \item historique de sélection et de réentraînement.
\end{itemize}

Les résultats négatifs, runs échoués et checkpoints non retenus
doivent être conservés dans les journaux.

% ============================================================
\section{Conclusion du protocole}
% ============================================================

EX12 teste une modification explicite d'AdaptiveRBF :
les centres ne sont plus seulement des unités locales optimisées par
la perte globale ; ils sont également évalués selon leur contribution,
guidés vers les erreurs restantes et organisés par un graphe de
relations entre leurs empreintes.

Le protocole sépare ces mécanismes afin de déterminer :
\begin{enumerate}
  \item si les centres peuvent être sélectionnés de façon fiable ;
  \item si les centres faibles peuvent devenir complémentaires ;
  \item si les empreintes capturent une structure utile ;
  \item si cette structure améliore les résultats au-delà d'un
  AdaptiveRBF classique à budget comparable.
\end{enumerate}

Aucune supériorité n'est présumée. La conclusion dépendra des résultats
confirmatoires, de la stabilité entre graines et des ablations.

\end{document}